% ═══════════════════════════════════════════════════════════════════════════════
% BC Sign and Minimum Analysis
% ═══════════════════════════════════════════════════════════════════════════════
% TITLE: Sign Structure of V'_lin(d) Under Various Boundary Conditions
% STATUS: [Der] - Derived from Green's function mathematics
% COMPANION: aside_frozen_brane_bc_v1/05_SIGN_AND_MINIMUM_ANALYSIS.md (original)
% ═══════════════════════════════════════════════════════════════════════════════

\documentclass[11pt,a4paper]{article}
\usepackage{amsmath,amssymb,amsthm}
\usepackage{booktabs}
\usepackage[margin=2.5cm]{geometry}

\newtheorem{theorem}{Theorem}
\newtheorem{lemma}{Lemma}

\title{Sign Structure of $V'_{\mathrm{lin}}(d)$ Under Various Boundary Conditions}
\author{EDC Derivation Library}
\date{January 2026}

\begin{document}
\maketitle

\begin{abstract}
We analyze the sign of the linearized vortex--vortex interaction potential derivative
$V'_{\mathrm{lin}}(d)$ under Neumann, Dirichlet, Robin, and ``frozen brane'' boundary
conditions. The conclusion is universal: $V'_{\mathrm{lin}}(d) > 0$ for all $d > 0$
regardless of BC choice. The minimum in the full potential arises from topological
core repulsion, not from boundary conditions.
\end{abstract}

\section{The Original No-Go Result}

\begin{lemma}[Monotonicity under Neumann BC]
For same-sign vortices under Neumann BC with uniform $\xi$-weight:
\begin{equation}
V'_{\mathrm{lin}}(d) = \frac{1}{d} + \frac{2}{\delta} \sum_{m=1}^{\infty}
\frac{1}{m} K_1\!\left(\frac{m\pi d}{\delta}\right) > 0
\quad \text{for all } d > 0
\end{equation}
\end{lemma}

\begin{theorem}[No Minimum in Linearized Potential]
The linearized potential $V_{\mathrm{lin}}(d)$ has no local minimum for same-sign vortices.
\end{theorem}

\begin{theorem}[Core Divergence]
$V_{\mathrm{core}}(d) \to +\infty$ as $d \to 0$ from topological winding constraints.
\end{theorem}

\begin{theorem}[Minimum Existence]
The full potential $V(d) = V_{\mathrm{core}}(d) + V_{\mathrm{lin}}(d)$ has a minimum
at some $d_0 > 0$ due to the balance of core repulsion and logarithmic growth.
\end{theorem}

\section{General Structure Under Arbitrary BC}

Under any BC with mode spectrum $\{\mu_m\}$, the linearized potential takes the form:
\begin{equation}
V_{\mathrm{lin}}(d) = \sum_m c_m \, G_m(d)
\end{equation}
where:
\begin{itemize}
    \item $c_m = (n_1 n_2) \times (W_m^2 / N_m) \times (\text{coupling factor})$
    \item $G_0(d) = -\frac{1}{2\pi} \ln(d/L)$ if zero mode exists
    \item $G_m(d) = \frac{1}{2\pi} K_0(\mu_m d)$ for massive modes
\end{itemize}

\subsection{Sign of Coefficients}

For same-sign vortices ($n_1 n_2 > 0$):
\begin{itemize}
    \item $c_m$ has definite sign (positive for repulsion convention)
    \item $W_m^2 \geq 0$ always (squared quantity)
    \item $N_m > 0$ always (integral of positive quantity)
\end{itemize}

\textbf{Conclusion:} No coefficient can be negative.

\subsection{Derivatives}

\textbf{Zero mode:}
\begin{equation}
\frac{d}{dd}\left[\ln d\right] = \frac{1}{d} > 0
\end{equation}

\textbf{Massive modes:}
\begin{equation}
\frac{d}{dd}\left[K_0(\mu d)\right] = -\mu K_1(\mu d) < 0
\end{equation}

But entering $V_{\mathrm{lin}}$ with negative overall coefficient:
\begin{equation}
V_{\mathrm{lin}} \propto -c_m K_0(\mu_m d)
\quad\Rightarrow\quad
V'_{\mathrm{lin}} \propto +c_m \mu_m K_1(\mu_m d) > 0
\end{equation}

\textbf{Every term in $V'_{\mathrm{lin}}$ is positive.}

\section{Explicit Calculation: Dirichlet BC}

\subsection{Mode Spectrum}
\begin{equation}
\mu_m = \frac{m\pi}{\delta}, \quad m = 1, 2, 3, \ldots \quad (\text{no zero mode})
\end{equation}

\subsection{Derivative}
\begin{equation}
V'_{\mathrm{lin}}(d) = n_1 n_2 \times \frac{8}{\pi\delta^2}
\sum_{k=0}^{\infty} \frac{1}{2k+1} K_1\!\left(\frac{(2k+1)\pi d}{\delta}\right)
\end{equation}

\textbf{Every term is positive for $n_1 n_2 > 0$.}

\section{Robin BC Analysis}

Robin BC: $\partial_\xi\phi|_0 = \kappa\phi|_0$, $\partial_\xi\phi|_\delta = -\kappa\phi|_\delta$

The transcendental equation:
\begin{equation}
\tan(\mu\delta) = \frac{2\kappa\mu}{\mu^2 - \kappa^2}
\end{equation}

In all cases:
\begin{equation}
V'_{\mathrm{lin}}(d) = \sum_m (\text{positive coefficient}) \times \mu_m \times K_1(\mu_m d) > 0
\end{equation}

\textbf{No regime of $\kappa$ creates $V' < 0$.}

\section{Frozen (Stiff) Limit}

``Frozen on brane'' = Dirichlet with matching boundary values: $\phi|_{\xi=0} = v$, $\phi|_{\xi=\delta} = v$.

The constant $v$ acts as a zero mode, giving:
\begin{equation}
V_{\mathrm{lin}}(d) \propto \ln(d) \quad\Rightarrow\quad V'_{\mathrm{lin}}(d) \propto \frac{1}{d} + [K_1 \text{ terms}] > 0
\end{equation}

\textbf{Still monotonically increasing.}

\section{Summary Table}
\label{sec:bc-sign-summary}

\begin{center}
\begin{tabular}{@{}lcccc@{}}
\toprule
BC Type & Zero Mode & $V_{\mathrm{lin}}(d \to 0)$ & $V_{\mathrm{lin}}(d \to \infty)$ & $V'_{\mathrm{lin}}(d)$ \\
\midrule
Neumann & Yes & $-\infty$ & $+\infty$ & $\mathbf{>0}$ \\
Dirichlet & No & finite & $+\infty$ & $\mathbf{>0}$ \\
Robin & No & finite & $+\infty$ & $\mathbf{>0}$ \\
Frozen & Yes$^*$ & $-\infty$ & $+\infty$ & $\mathbf{>0}$ \\
\bottomrule
\end{tabular}
\end{center}

\section{Conclusion}

\textbf{$V'_{\mathrm{lin}}(d) > 0$ for all BC scenarios analyzed.}

The brane-frozen regime (Dirichlet with stiff pinning) does NOT:
\begin{itemize}
    \item Create an attractive window
    \item Change the monotonicity of $V_{\mathrm{lin}}$
    \item Move the minimum location significantly
\end{itemize}

\textbf{The minimum in $V(d)$ comes from core overlap (topology), not from BC.}

\subsection{What BC Actually Contribute}

\textbf{Positive contributions:}
\begin{enumerate}
    \item \textbf{Scale $\delta$:} BC set the brane thickness scale
    \item \textbf{Mode spectrum:} BC determine eigenvalues $\mu_m$
    \item \textbf{Screening:} Higher modes provide exponential screening at $d > \delta$
\end{enumerate}

\textbf{What BC do NOT contribute:}
\begin{enumerate}
    \item \textbf{Attraction:} No BC creates $V' < 0$
    \item \textbf{Minimum mechanism:} Comes from topology, not BC
    \item \textbf{Sign reversal:} Impossible with standard Green's functions
\end{enumerate}

\end{document}
